\documentclass[11pt,a4paper]{article}
\usepackage[T1]{fontenc}
\usepackage[utf8]{inputenc}
\usepackage{amsmath,amssymb,amsthm}
\usepackage[margin=1in]{geometry}
\usepackage[colorlinks=true,linkcolor=blue,urlcolor=blue]{hyperref}
\theoremstyle{plain}
\newtheorem{theorem}{Theorem}
\newtheorem{lemma}[theorem]{Lemma}
\newtheorem{proposition}[theorem]{Proposition}
\newtheorem{corollary}[theorem]{Corollary}
\theoremstyle{definition}
\newtheorem{remark}[theorem]{Remark}

\title{The empirical recurrence for OEIS A184141, proved by transfer matrix}
\author{Adrian Perez Fontelles\\ \small Independent researcher}
\date{13 September 2026}
\begin{document}
\maketitle

\begin{abstract}
OEIS A184141 counts arrays over $\{0,\dots,1\}$ whose rows, columns, diagonals or
antidiagonals are required to be monotone or unimodal, and carries an empirical recurrence of
order $64$ contributed by R.~H.~Hardin. A monotone condition is local, but a unimodal one
is not: a sequence is unimodal exactly when it never rises again after it has fallen, so the
state must remember which sequences have already fallen. With those bits carried along, $a(n)$
is a walk count on $S=2048$ vertices and the recurrence follows from a finite exact
computation.
\end{abstract}

\noindent\small 2020 Mathematics Subject Classification. 05A15, 05A16, 15B36.\normalsize

\section{The sequence and the conjecture}

OEIS A184141 is ``Number of n X 6 binary arrays with rows and columns in nondecreasing order.''. It has offset $1$ and begins
\[
7, 63, 785, 11819, 183010, 2625117, 33345183, 371484319,\ \dots
\]
The entry carries the following comment:
\begin{quote}\small
Empirical: a(n)=64*a(n-1)-2016*a(n-2)+41664*a(n-3)-635376*a(n-4)+7624512*a(n-5)-74974368*a(n-6)+621216192*a(n-7)-4426165368*a(n-8)+27540584512*a(n-9)-151473214816*a(n-10)+743595781824*a(n-11)-3284214703056*a(n-12)+13136858812224*a(n-13)-47855699958816*a(n-14)+159518999862720*a(n-15)-488526937079580*a(n-16)+1379370175283520*a(n-17)-3601688791018080*a(n-18)+8719878125622720*a(n-19)-19619725782651120*a(n-20)+41107996877935680*a(n-21)-80347448443237920*a(n-22)+146721427591999680*a(n-23)-250649105469666120*a(n-24)+401038568751465792*a(n-25)-601557853127198688*a(n-26)+846636978475316672*a(n-27)-1118770292985239888*a(n-28)+1388818294740297792*a(n-29)-1620288010530347424*a(n-30)+1777090076065542336*a(n-31)-1832624140942590534*a(n-32)+1777090076065542336*a(n-33)-1620288010530347424*a(n-34)+1388818294740297792*a(n-35)-1118770292985239888*a(n-36)+846636978475316672*a(n-37)-601557853127198688*a(n-38)+401038568751465792*a(n-39)-250649105469666120*a(n-40)+146721427591999680*a(n-41)-80347448443237920*a(n-42)+41107996877935680*a(n-43)-19619725782651120*a(n-44)+8719878125622720*a(n-45)-3601688791018080*a(n-46)+1379370175283520*a(n-47)-488526937079580*a(n-48)+159518999862720*a(n-49)-47855699958816*a(n-50)+13136858812224*a(n-51)-3284214703056*a(n-52)+743595781824*a(n-53)-151473214816*a(n-54)+27540584512*a(n-55)-4426165368*a(n-56)+621216192*a(n-57)-74974368*a(n-58)+7624512*a(n-59)-635376*a(n-60)+41664*a(n-61)-2016*a(n-62)+64*a(n-63)-a(n-64) (=polynomial degree 63).
\end{quote}
As of the ``Last modified'' line on the live entry (Aug 13 2026 15:26:08, revision 13) the statement is
still recorded as empirical, and nothing on the entry records it as settled either way.

Written out, the claim is
\begin{equation}\label{eq:conj}
a(n)\;=\;64\,a(n-1) - 2016\,a(n-2) + 41664\,a(n-3) - 635376\,a(n-4) + 7624512\,a(n-5) - 74974368\,a(n-6) + 621216192\,a(n-7) - 4426165368\,a(n-8) + 27540584512\,a(n-9) - 151473214816\,a(n-10) + 743595781824\,a(n-11) - 3284214703056\,a(n-12) + 13136858812224\,a(n-13) - 47855699958816\,a(n-14) + 159518999862720\,a(n-15) - 488526937079580\,a(n-16) + 1379370175283520\,a(n-17) - 3601688791018080\,a(n-18) + 8719878125622720\,a(n-19) - 19619725782651120\,a(n-20) + 41107996877935680\,a(n-21) - 80347448443237920\,a(n-22) + 146721427591999680\,a(n-23) - 250649105469666120\,a(n-24) + 401038568751465792\,a(n-25) - 601557853127198688\,a(n-26) + 846636978475316672\,a(n-27) - 1118770292985239888\,a(n-28) + 1388818294740297792\,a(n-29) - 1620288010530347424\,a(n-30) + 1777090076065542336\,a(n-31) - 1832624140942590534\,a(n-32) + 1777090076065542336\,a(n-33) - 1620288010530347424\,a(n-34) + 1388818294740297792\,a(n-35) - 1118770292985239888\,a(n-36) + 846636978475316672\,a(n-37) - 601557853127198688\,a(n-38) + 401038568751465792\,a(n-39) - 250649105469666120\,a(n-40) + 146721427591999680\,a(n-41) - 80347448443237920\,a(n-42) + 41107996877935680\,a(n-43) - 19619725782651120\,a(n-44) + 8719878125622720\,a(n-45) - 3601688791018080\,a(n-46) + 1379370175283520\,a(n-47) - 488526937079580\,a(n-48) + 159518999862720\,a(n-49) - 47855699958816\,a(n-50) + 13136858812224\,a(n-51) - 3284214703056\,a(n-52) + 743595781824\,a(n-53) - 151473214816\,a(n-54) + 27540584512\,a(n-55) - 4426165368\,a(n-56) + 621216192\,a(n-57) - 74974368\,a(n-58) + 7624512\,a(n-59) - 635376\,a(n-60) + 41664\,a(n-61) - 2016\,a(n-62) + 64\,a(n-63) - a(n-64)
\end{equation}
for all $n>64$.

\section{The conditions}

The arrays have $6$ columns and $L=n$ rows. Written in (row, column)
coordinates, the entry asks:
\begin{itemize}
\item the rows themselves, read as vectors, are in lexicographically nondecreasing order;
\item the columns themselves, read as vectors, are in lexicographically nondecreasing order;
\end{itemize}
Here ``unimodal'' means nondecreasing and then nonincreasing, that is: no rise after a fall.

A lexicographic condition between adjacent columns is decided
at the first row where they differ, so the state carries, for each adjacent pair of
columns, one bit: still tied, or already settled in the right order. A row that would
settle a pair the wrong way has no outgoing edge.

\begin{lemma}\label{lem:walk}
Let $\mathcal V$ be the set of pairs (a row satisfying every within-row condition, a
bit vector), $|\mathcal V|=S=2048$; let $(r,f)\to(s,f')$ be an edge when placing $s$ after $r$
breaks none of the across-row conditions and $f'$ is $f$ updated by that step; let
$\mathbf{v}$ mark the vertices whose bit vector is the one their row alone produces.
Then
\[
a(n)\;=\;\,\mathbf{v}^{\!\top}N^{\,L-1}\mathbf{1},\qquad L=n.
\]
\end{lemma}

\begin{proof}
A walk of length $L-1$ is a sequence of $L$ rows, that is, an array, and its bit component
is determined by them. Every within-row condition is imposed on each vertex; every
across-row condition compares two consecutive rows and is imposed on each edge; and the
unimodal and lexicographic bits are, by construction, the record of what the rows laid down
so far have already decided, so a violation removes the corresponding edge. Hence the walks
are exactly the admissible arrays. Nothing constrains the last row beyond what has already
been imposed, so the walk may end anywhere. The state carries a row, not a boundary, so
$L$ rows make a walk of length $L-1$.
\end{proof}

\section{The criterion, and the computation}

Let $q(t)=t^{64}-\sum_i c_i\,t^{64-i}$ be the characteristic polynomial of
\eqref{eq:conj} and $G=N$.

\begin{lemma}\label{lem:crit}
With $n_0$ the least index for which $L\ge1$, $o=1n_0+0-1$ and
$h_j=G^{j}N^{o}\mathbf{1}$, we have for $j\ge64$
\[
a(n_0+j)-\sum_i c_i\,a(n_0+j-i)=\,\mathbf{v}^{\!\top}G^{\,j-64}w,
\qquad w=h_{64}-\sum_i c_ih_{64-i} .
\]
Hence \eqref{eq:conj} holds from that point on if and only if
$\mathbf{v}^{\!\top}G^{m}w=0$ for all $m\ge0$, and $m<S$ suffices.
\end{lemma}

\begin{proof}
Substitute Lemma~\ref{lem:walk} and factor. By Cayley--Hamilton $G^{S}$ is an integer
combination of $I,G,\dots,G^{S-1}$, so every later value repeats that combination of
earlier ones.
\end{proof}

\begin{theorem}
The recurrence \eqref{eq:conj} holds for every $n>64$, which is the statement of
Section~1.
\end{theorem}

\begin{proof}
The digraph was built directly from the conditions listed above. The vector $w$ was formed by
$64$ applications of $G$ in exact integer arithmetic and $\mathbf{v}^{\!\top}G^{m}w$
evaluated for $m=0,1,\dots$, again exactly. Every one of those integers vanishes from the
index corresponding to $n=64+1$ onwards, and the run continued until $S+1$ consecutive
zeros had been seen.
\end{proof}

\section{Verification}

Three checks, each independent of the algebra above.

First, the walk counts of Lemma~\ref{lem:walk} were compared against the entry's DATA at
every one of the $21$ published indices, with no shift fitted. The values agree exactly.
This is what fixes the reading of the wording: ``rows and columns in nondecreasing order''
means the rows and columns are sorted as vectors, not that each is nondecreasing along
itself, and the two readings already differ at the second published term --- $86$ against
$50$ for $2\times2$ arrays over $\{0,1,2,3\}$. The DATA decides it.

Second, the recurrence \eqref{eq:conj} was evaluated on those published terms in exact
integer arithmetic, with no matrices involved, and vanishes wherever it applies.

Third, the annihilation test of Lemma~\ref{lem:crit} was carried to the full
Cayley--Hamilton bound in exact integer arithmetic rather than to a sampled prefix.

\begin{thebibliography}{9}
\bibitem{oeis} The OEIS Foundation, \emph{The On-Line Encyclopedia of Integer
Sequences}, \url{https://oeis.org}, 2026. Sequence A184141.
\bibitem{stanley} R.~P.~Stanley, \emph{Enumerative Combinatorics}, Volume 1, 2nd ed.,
Cambridge University Press, 2011. (Transfer-matrix method, Section 4.7.)
\bibitem{fs} P.~Flajolet and R.~Sedgewick, \emph{Analytic Combinatorics}, Cambridge
University Press, 2009.
\bibitem{horn} R.~A.~Horn and C.~R.~Johnson, \emph{Matrix Analysis}, 2nd ed., Cambridge
University Press, 2013.
\end{thebibliography}

\end{document}
